% ================= CAPÍTULO 5: DETERMINANTES =================
\chapter{Determinantes}
\salio{9}{«$\det A=k$, calcular $\det B$» o escalerizar; más 8 de expresiones matriciales}

\section{Definición y cálculo}
\begin{sello}{DEFINICIÓN (DESARROLLO POR UNA FILA)}
$\det(a)=a$. \quad $\begin{vmatrix}a&b\\c&d\end{vmatrix}=ad-bc$.\\[3pt]
Para $n\times n$: $\det A=\displaystyle\sum_{j=1}^n(-1)^{i+j}a_{ij}\det A_{ij}$, donde $A_{ij}$ es la matriz que queda al \textbf{borrar la fila $i$ y la columna $j$}.\\[2pt]
\textbf{Teorema:} da lo mismo para cualquier fila $i$, y también desarrollando por cualquier columna.
\end{sello}
\begin{memo}{SIGNOS Y SARRUS}
Los signos $(-1)^{i+j}$ forman un tablero de ajedrez: $\begin{smallmatrix}+&-&+\\-&+&-\\+&-&+\end{smallmatrix}$.\\[2pt]
Regla de Sarrus ($3\times3$, \textbf{solo} $3\times3$):\\
$\begin{vmatrix}a&b&c\\d&e&f\\g&h&i\end{vmatrix}=aei+bfg+cdh-ceg-afh-bdi$.
\end{memo}
\begin{metodo}{EL CAMINO RÁPIDO PARA MATRICES GRANDES}
1) Desarrollá por la fila o columna con \textbf{más ceros}.\\
2) Si no hay ceros, fabricalos: $F_i\to F_i+\lambda F_j$ \textbf{no cambia} el determinante.\\
3) Llevá a triangular: el determinante es el \textbf{producto de la diagonal}.\\[2pt]
1S 2019: una $6\times6$. Restando la fila 1 a todas las demás queda triangular con diagonal $1,-2,-2,-2,1,1$: $\det=-8$.
\end{metodo}

\section{Propiedades}
\begin{sello}{CÓMO REACCIONA A LAS OPERACIONES}
1) \textbf{Intercambiar} dos filas: cambia el signo.\\
2) \textbf{Multiplicar una fila} por $\lambda$: el determinante queda multiplicado por $\lambda$.\\
3) \textbf{Sumarle a una fila un múltiplo de otra}: no cambia.\\
4) $\det(A^t)=\det A$: todo lo anterior vale igual para \textbf{columnas}.\\
5) \textbf{Lineal en cada fila:} si una fila es suma, el determinante se parte en dos:
\[\begin{vmatrix}u+v\\ F_2\\ F_3\end{vmatrix}=\begin{vmatrix}u\\ F_2\\ F_3\end{vmatrix}+\begin{vmatrix}v\\ F_2\\ F_3\end{vmatrix}.\]
\end{sello}
\begin{sello}{CONSECUENCIAS}
Fila (o columna) de ceros, dos filas iguales, o dos filas proporcionales $\Rightarrow\det=0$.\\
Triangular (superior o inferior) $\Rightarrow\det$ = producto de la diagonal. En particular $\det I=1$.\\
$\boldsymbol{\det(\lambda A)=\lambda^n\det A}$ para $A$ de $n\times n$ ($\lambda$ sale una vez \emph{por fila}).
\end{sello}
\begin{trampa}{LAS DOS QUE MÁS SE EQUIVOCAN}
$\det(2A)=2\det A$ es \textbf{falso}: es $2^n\det A$. (1S 2019, 2S 2018, 2S 2019.)\\
$\det(A+B)=\det A+\det B$ es \textbf{falso}. La linealidad es \emph{de a una fila}, no de toda la matriz.\\
Y «$\det A=0\Rightarrow A$ tiene una fila o columna de ceros» también es falso: $\begin{pmatrix}1&2\\2&4\end{pmatrix}$.
\end{trampa}

\section{El determinante del producto}
\begin{sello}{TEOREMA}
$\boldsymbol{\det(AB)=\det A\cdot\det B}$ \quad ($A,B$ de $n\times n$).\\[3pt]
De ahí: \ $\det(A^k)=(\det A)^k$ \quad $\det(A^{-1})=\dfrac1{\det A}$ \quad $\det(P^{-1}AP)=\det A$.\\[3pt]
$A$ es invertible $\iff\det A\neq0$.
\end{sello}
\begin{demo}
$\det(A^{-1})$: de $AA^{-1}=I$ sale $\det A\cdot\det A^{-1}=\det I=1$.\\
Invertible $\Rightarrow\det\neq0$: por lo anterior, $\det A\cdot\det A^{-1}=1$.\\
$\det\neq0\Rightarrow$ invertible: al escalerizar, cada operación multiplica el determinante por un número no nulo. Si $\det A\neq0$, la escalonada tiene determinante no nulo, así que no tiene filas de ceros: rango $n$, invertible. $\blacksquare$
\end{demo}
\begin{metodo}{DETERMINANTE DE UNA EXPRESIÓN MATRICIAL}
1) \textbf{Factorizá} la expresión en un producto (respetando izquierda y derecha).\\
2) Aplicá $\det(XY)=\det X\det Y$, $\det(\lambda X)=\lambda^n\det X$, $\det X^{-1}=1/\det X$, $\det X^t=\det X$.\\
3) Despejá el número que te piden.
\end{metodo}
\begin{memo}{FACTORIZACIONES QUE SALIERON}
$A^2+AB-BA-B^2=A(A+B)-B(A+B)=(A-B)(A+B)$ \quad (2S 2017).\\
$A^2+2AB-2BA-4B^2=(A-2B)(A+2B)$ \quad (1S 2023).\\
$B^3A^3B^{-1}+B^4A^2B^{-1}=B^3(A^3+BA^2)B^{-1}=B^3(A+B)A^2B^{-1}$ \quad (1S 2023 vesp.).\\
Fijate que $A^3+BA^2=(A+B)A^2$: el $A^2$ sale por la \textbf{derecha}.
\end{memo}
\begin{memo}{ECUACIONES MATRICIALES}
2S 2016 ($2\times2$): $3A^{-1}B^2=2I$, $\det B=2$. \ $\det$: $3^2\cdot\frac1{\det A}\cdot4=2^2\Rightarrow\det A=9$.\\
1S 2018 ($3\times3$): $B^{-1}A(B^t)^2=-2I$, $\det B=2$. \ $\frac12\det A\cdot4=(-2)^3\Rightarrow\det A=-4$.\\
2S 2019 ($5\times5$): $|A|=2\Rightarrow|2A^{-1}|=2^5\cdot\frac12=16$. «Vale 1» es falso.
\end{memo}

\section{Casos especiales}
\salio{5}{antisimétricas, nilpotentes, idempotentes}
\begin{sello}{LO QUE SE SABE DEL DETERMINANTE SIN CALCULARLO}
\textbf{Antisimétrica con $n$ impar:} $\det A=0$.\\
\textbf{Nilpotente} ($A^k=O$): $\det A=0$.\\
\textbf{Idempotente} ($A^2=A$): $\det A\in\{0,1\}$.\\
\textbf{Involutiva} ($A^2=I$): $\det A=\pm1$.\\
$A^3=A$ con $\det A\neq0$: $\det A=\pm1$ y $A=A^{-1}$.
\end{sello}
\begin{demo}
Antisimétrica: $\det A=\det A^t=\det(-A)=(-1)^n\det A=-\det A$, así que $\det A=0$.\\
Nilpotente: $(\det A)^k=\det(A^k)=\det O=0$.\\
Idempotente: $(\det A)^2=\det A$, así que $\det A(\det A-1)=0$.\\
Involutiva: $(\det A)^2=1$. \ $A^3=A$: $(\det A)^3=\det A\neq0\Rightarrow(\det A)^2=1$; y multiplicando $A^3=A$ por $A^{-1}$: $A^2=I$. $\blacksquare$
\end{demo}
\begin{memo}{1S 2018 · ANTISIMÉTRICA NO NULA, $n$ IMPAR}
$A$ no es invertible ($\det=0$), y $A^tA$ tampoco: $\det(A^tA)=(\det A)^2=0$. Respuesta: ninguna de las dos.
\end{memo}
\begin{sello}{BLOQUES Y VANDERMONDE}
$\det\begin{pmatrix}A&B\\O&C\end{pmatrix}=\det A\cdot\det C$ \quad (práctico 5, ej. 6).\\[3pt]
$\begin{vmatrix}1&1&1\\a&b&c\\a^2&b^2&c^2\end{vmatrix}=(b-a)(c-a)(c-b)$: es invertible $\iff a,b,c$ son distintos.\\
2S 2022: $\rg(V)=1\iff a=b=c$.
\end{sello}

\section{El método estrella: «$\det A=k$, ¿cuánto vale $\det B$?»}
\salio{5}{1S 2017, 2S 2017, 2S 2018, 2S 2024, 1S 2025}
\begin{metodo}{LLEVAR $B$ HASTA $A$ ANOTANDO CADA FACTOR}
1) Sacá los factores comunes de cada fila y de cada columna (cada uno multiplica).\\
2) Deshacé las sumas: $F_i\to F_i-\lambda F_j$ (no cambia nada). Si una fila es suma de dos cosas que no se cancelan, partila (linealidad).\\
3) Reordená filas y columnas hasta llegar a $A$: \textbf{cada intercambio cambia el signo}. Una permutación cíclica de 3 filas son 2 intercambios: signo $+$.\\
4) Multiplicá todos los factores por $k$.
\end{metodo}
\begin{memo}{PRÁCTICO 5 · EJ. 2 (CON $\det A=5$)}
$(d,e,f;\ g,h,j;\ a,b,c)$: cíclica, dos intercambios $\to5$.\\
$(-a,-b,-c;\ 2d,2e,2f;\ -g,-h,-j)$: $(-1)(2)(-1)=2\ \to10$.\\
$(a,b,c;\ d-3a,\dots;\ 2g,2h,2j)$: la resta no cambia, el 2 sí $\to10$.\\
$(a+5c,3b,c;\ \dots;\ 2g+10j,6h,2j)$: sale 2 de la fila 3, sale 3 de la columna 2, y $C_1\to C_1-5C_3$ $\to30$.
\end{memo}

\section{Lo que significa: área y volumen}
\begin{tip}{LA IMAGEN QUE LO EXPLICA TODO}
$|\det A|$ es el \textbf{factor por el que la transformación $X\mapsto AX$ multiplica las áreas} (en $3\times3$, los volúmenes). El cuadradito unitario se convierte en el paralelogramo formado por las columnas de $A$, y su área es $|\det A|$.\\[2pt]
\textbf{El signo} dice si la transformación da vuelta el plano como un espejo (negativo) o no (positivo).\\[2pt]
\textbf{$\det A=0$} quiere decir que el plano se \textbf{aplasta} en una recta o en un punto: se pierde información y no hay forma de volver atrás. Por eso $\det A=0\iff A$ no es invertible.\\[2pt]
$\det(AB)=\det A\det B$: si una máquina multiplica las áreas por 3 y la otra por 2, hacer las dos las multiplica por 6.
\end{tip}
\begin{center}
\begin{tikzpicture}[scale=0.75,>=Stealth]
\fill[acero!25] (0,0) rectangle (1,1); \draw[acero] (0,0) rectangle (1,1);
\node[font=\scriptsize] at (0.5,-0.7) {área 1};
\draw[->,thick,suave] (1.4,0.5)--(2.2,0.5);
\begin{scope}[xshift=2.7cm]
\fill[rojo!20] (0,0)--(3,0)--(4,2)--(1,2)--cycle; \draw[rojo] (0,0)--(3,0)--(4,2)--(1,2)--cycle;
\draw[->,thick,rojo] (0,0)--(3,0); \draw[->,thick,teal] (0,0)--(1,2);
\node[font=\scriptsize,align=center] at (2,-0.9) {$\begin{pmatrix}3&1\\0&2\end{pmatrix}$\\ área $6=\det$};
\end{scope}
\begin{scope}[xshift=8.2cm]
\draw[rojo,dashed] (-0.3,-0.6)--(1.4,2.8);
\draw[->,thick,teal] (0,0)--(1.2,2.4); \draw[->,thick,rojo] (0,0)--(0.6,1.2);
\node[font=\scriptsize,align=center] at (0.6,-0.9) {$\begin{pmatrix}1&2\\2&4\end{pmatrix}$\\ $\det=0$: se aplasta};
\end{scope}
\end{tikzpicture}
\end{center}

\section{Con parámetro: ¿para qué valores es invertible?}
\begin{metodo}{PRÁCTICO 5 · EJ. 3}
1) Calculá $\det A$ como polinomio en $k$ (desarrollo o escalerizando con cuidado).\\
2) \textbf{Factorizá.} $A$ es invertible exactamente cuando ese polinomio no se anula.\\[2pt]
Ej. 3a: $\begin{vmatrix}k&-k&3\\0&k+1&1\\k&-8&k-1\end{vmatrix}=k(k-2)^2$: invertible $\iff k\neq0$ y $k\neq2$.\\
2S 2017: $\det=-2k^2+8k+24=-2(k+2)(k-6)$: positivo entre $-2$ y $6$, negativo afuera.
\end{metodo}
\begin{sello}{INVERSA CON LA ADJUNTA Y REGLA DE CRAMER}
$A^{-1}=\dfrac1{\det A}\,\mathrm{Adj}(A)$, con $\mathrm{Adj}(A)_{ij}=(-1)^{i+j}\det A_{ji}$ (la traspuesta de la matriz de cofactores). En $2\times2$ es la fórmula del capítulo 4.\\[3pt]
\textbf{Cramer:} si $\det A\neq0$, la solución de $AX=b$ es $x_j=\dfrac{\det A_j}{\det A}$, donde $A_j$ es $A$ con la columna $j$ cambiada por $b$. Sirve cuando piden \textbf{una sola} incógnita.
\end{sello}
\begin{memo}{PRÁCTICO 5 · EJ. 5}
$d_n$: la matriz de unos con ceros en la diagonal salvo el primer lugar. Restando la fila 1 a cada una de las otras, la fila $i$ queda $-e_i$. Queda triangular con diagonal $1,-1,\dots,-1$: $\ d_n=(-1)^{n-1}$.
\end{memo}
